% !TEX program = pdflatex
\documentclass[aspectratio=169, 11pt]{beamer}

% ── theme & colors ──────────────────────────────────────────────────────────
\usetheme{Madrid}
\usecolortheme{default}
\definecolor{darkblue}{RGB}{44,62,80}
\definecolor{accentred}{RGB}{231,76,60}
\definecolor{accentgreen}{RGB}{39,174,96}
\definecolor{accentblue}{RGB}{52,152,219}
\definecolor{lightgray}{RGB}{245,245,245}

\setbeamercolor{structure}{fg=darkblue}
\setbeamercolor{title}{fg=white, bg=darkblue}
\setbeamercolor{frametitle}{fg=darkblue}
\setbeamercolor{block title}{bg=darkblue, fg=white}
\setbeamercolor{block body}{bg=lightgray}
\setbeamercolor{alerted text}{fg=accentred}
\setbeamertemplate{navigation symbols}{}
\setbeamertemplate{footline}{%
  \leavevmode\hbox{%
    \begin{beamercolorbox}[wd=.333\paperwidth,ht=2.5ex,dp=1ex,center]{author in head/foot}%
      \usebeamerfont{author in head/foot}\insertshortauthor
    \end{beamercolorbox}%
    \begin{beamercolorbox}[wd=.334\paperwidth,ht=2.5ex,dp=1ex,center]{title in head/foot}%
      \usebeamerfont{title in head/foot}\insertshorttitle
    \end{beamercolorbox}%
    \begin{beamercolorbox}[wd=.333\paperwidth,ht=2.5ex,dp=1ex,right]{date in head/foot}%
      \usebeamerfont{date in head/foot}\insertframenumber{}/\inserttotalframenumber\hspace*{2ex}
    \end{beamercolorbox}%
  }%
}

% ── packages ────────────────────────────────────────────────────────────────
\usepackage{amsmath, amssymb}
\usepackage{booktabs}
\usepackage{graphicx}
\usepackage{tikz}
\usepackage{multicol}
\usepackage{adjustbox}

% ── metadata ────────────────────────────────────────────────────────────────
\title[Japan CPI Assemblage Regression]{%
  Assemblage Regression for\\Japan CPI Inflation Forecasting}
\subtitle{Component-Space and Rank-Space Approaches}
\author[F.\ Yuan]{Fangyuan}
\institute[MSFE]{Master of Science in Financial Engineering}
\date{\today}

% ════════════════════════════════════════════════════════════════════════════
\begin{document}

% ── Slide 1: Title ──────────────────────────────────────────────────────────
\begin{frame}[plain]
  \titlepage
\end{frame}

% ── Slide 2: Outline ───────────────────────────────────────────────────────
\begin{frame}{Outline}
  \tableofcontents
\end{frame}

% ════════════════════════════════════════════════════════════════════════════
\section{Motivation \& Data}
% ════════════════════════════════════════════════════════════════════════════

% ── Slide 3: Motivation ────────────────────────────────────────────────────
\begin{frame}{Motivation}
  \begin{itemize}
    \item Japan experienced decades of deflation (1990s--2010s), followed by a sharp
          inflation surge in 2022--2024.
    \item Standard ``core'' measures (ex fresh food, ex food \& energy) are
          \alert{fixed-weight subsets} --- they cannot adapt to changing inflation dynamics.
    \item \textbf{Research question:} Can we learn \emph{data-driven} component weights
          that are more predictive of future headline inflation than
          traditional core measures?
  \end{itemize}
  \vfill
  \begin{block}{Approach: Assemblage Regression}
    Constrained ridge regression that finds optimal component weights, shrinking
    toward official CPI basket weights as an informative prior.
  \end{block}
\end{frame}

% ── Slide 4: Data Overview ─────────────────────────────────────────────────
\begin{frame}{Data: Japan CPI Structure}
  \textbf{Source:} Japan Statistics Bureau, 2020 base year.\\[6pt]
  \begin{columns}[T]
    \begin{column}{0.48\textwidth}
      \textbf{Hierarchy:}
      \begin{itemize}
        \item \textbf{Level 1:} 10 major groups
        \item \textbf{Level 2:} 48 mid-level components
        \item \textbf{Level 3:} 744 granular items
      \end{itemize}
      \vspace{6pt}
      \textbf{Time span:} 1970--2026 (monthly)\\
      \textbf{Regression window:} 1990--2026
    \end{column}
    \begin{column}{0.48\textwidth}
      \textbf{Column classification (Level 2):}
      \begin{itemize}
        \item 6 \emph{composite} measures (headline variants)
        \item 4 \emph{special} cross-cutting aggregates
        \item \alert{46 active components} (non-overlapping)
        \item 1 \emph{isolated:} Rent (18.3\%)
      \end{itemize}
      \vspace{6pt}
      Active components sum to \textbf{81.7\%} of the basket.
    \end{column}
  \end{columns}
\end{frame}

% ── Slide 5: Rent Isolation ────────────────────────────────────────────────
\begin{frame}{Key Data Insight: Rent Isolation}
  \begin{columns}[T]
    \begin{column}{0.55\textwidth}
      \textbf{Rent = 18.3\% of the CPI basket}, but:
      \begin{itemize}
        \item $\sim$85\% is \emph{imputed rent} (\begin{CJK}{UTF8}{min}持家の帰属家賃\end{CJK}) --- an administrative estimate for owner-occupiers.
        \item Market rent: only 2.5\% of basket.
        \item $\sigma(\text{3m/3m}) = 1.09$ --- \alert{lowest} of all 47 components.
        \item Acts as a \emph{weight sink} in constrained regression: absorbs weight without contributing signal.
      \end{itemize}
    \end{column}
    \begin{column}{0.42\textwidth}
      \begin{block}{Treatment}
        \begin{itemize}
          \item Included in feature set for completeness.
          \item Analyzed separately in all charts.
          \item No separate market-rent series available.
        \end{itemize}
      \end{block}
    \end{column}
  \end{columns}
\end{frame}

% ── Slide 6: Basket Weights ───────────────────────────────────────────────
\begin{frame}{CPI Basket Weights (2020 Base)}
  \begin{columns}[T]
    \begin{column}{0.48\textwidth}
      \small
      \textbf{Top 10 components by basket share:}
      \begin{enumerate}
        \item Meals outside the home \hfill 7.16\%
        \item Private transportation \hfill 5.52\%
        \item Recreational services \hfill 4.17\%
        \item Medical services \hfill 3.32\%
        \item Communication \hfill 3.26\%
        \item Cooked food \hfill 2.92\%
        \item Electric-ity \hfill 2.85\%
        \item Clothes \hfill 2.35\%
        \item Tutorial fees \hfill 2.17\%
        \item Public transportation \hfill 2.15\%
      \end{enumerate}
    \end{column}
    \begin{column}{0.48\textwidth}
      \small
      \textbf{12 sector groups:}\\[4pt]
      \begin{tabular}{ll}
        \toprule
        Group & Share \\
        \midrule
        Food (other) & $\sim$23\% \\
        Fresh food & $\sim$5\% \\
        Energy & $\sim$5\% \\
        Transport & $\sim$8\% \\
        Recreation & $\sim$9\% \\
        Healthcare & $\sim$5\% \\
        Clothing & $\sim$4\% \\
        Education & $\sim$4\% \\
        Housing (other) & $\sim$3\% \\
        Household goods & $\sim$4\% \\
        Communication & $\sim$3\% \\
        Other & $\sim$9\% \\
        \bottomrule
      \end{tabular}
    \end{column}
  \end{columns}
\end{frame}

% ── Slide 7: Inflation Dynamics ────────────────────────────────────────────
\begin{frame}{Japan CPI: Historical Inflation Dynamics}
  \begin{columns}[T]
    \begin{column}{0.55\textwidth}
      \textbf{Three eras:}
      \begin{enumerate}
        \item \textbf{Lost decades} (1990--2012): Persistent deflation, headline
              CPI mostly negative YoY.
        \item \textbf{Abenomics} (2013--2021): BOJ targets 2\%; achieved only
              briefly.  Core measures hovered near 0--1\%.
        \item \textbf{2022 surge}: Global supply shocks + weak yen pushed
              headline above 4\%.  Energy and food dominated.
      \end{enumerate}
    \end{column}
    \begin{column}{0.42\textwidth}
      \textbf{Key observation:}\\[4pt]
      Published core measures (ex fresh food, ex food \& energy) diverged
      sharply from headline during the 2022 surge --- motivating the search
      for adaptive component weights.
    \end{column}
  \end{columns}
  \vfill
  \centering\small
  \textit{See: fig1\_core\_measures\_yoy.png, fig2\_recent\_3m3m.png}
\end{frame}

% ════════════════════════════════════════════════════════════════════════════
\section{Methodology}
% ════════════════════════════════════════════════════════════════════════════

% ── Slide 8: Framework Overview ────────────────────────────────────────────
\begin{frame}{Assemblage Regression Framework}
  \textbf{Growth rates:} 3-month-over-3-month annualized:
  \[
    g_{i,t} = \left(\frac{\text{CPI}_{i,t}}{\text{CPI}_{i,t-3}} - 1\right) \times 4 \times 100
  \]

  \textbf{Target:} 12-month forward average of headline 3m/3m growth:
  \[
    y_t = \frac{1}{12}\sum_{h=1}^{12} g_{\text{headline},\, t+h}
  \]

  \textbf{Evaluation protocol:}
  \begin{itemize}
    \item \textbf{In-sample:} Full data, $\lambda$ selected by 10-fold expanding-window CV.
    \item \textbf{Out-of-sample:} Rolling 1-step predictions.
          $\lambda$ fixed from first 120 months (no look-ahead).
    \item \textbf{Windows:} Expanding (all history) and rolling 20-year.
  \end{itemize}
\end{frame}

% ── Slide 9: Albacorecomps ─────────────────────────────────────────────────
\begin{frame}{Model 1: Albacorecomps (Component Space)}
  \begin{block}{Optimization Problem}
    \[
      \min_{\mathbf{w}} \;\; \frac{1}{T}\sum_{t=1}^{T}\bigl(y_t - \mathbf{x}_t'\mathbf{w}\bigr)^2
      \;+\; \lambda \|\mathbf{w} - \mathbf{w}_{\text{prior}}\|^2
    \]
    \vspace{-6pt}
    \[
      \text{s.t.} \quad w_i \geq 0, \quad \sum_i w_i = 1
    \]
  \end{block}
  \vspace{4pt}
  \begin{columns}[T]
    \begin{column}{0.48\textwidth}
      \textbf{Key design choices:}
      \begin{itemize}
        \item $\mathbf{X}$: $T \times K$ matrix of component 3m/3m growth rates
              ($K = 47$ for Level~2, $K = 9$ for Level~1).
        \item $\mathbf{w}_{\text{prior}}$: Official basket weights
              (renormalized to sum to 1).
        \item Non-negativity + sum-to-one: weights are interpretable
              as a ``core measure.''
      \end{itemize}
    \end{column}
    \begin{column}{0.48\textwidth}
      \textbf{Penalty structure:}
      \begin{itemize}
        \item Flat penalty: every component penalized equally for
              deviating from prior.
        \item $\lambda \to 0$: unconstrained least squares.
        \item $\lambda \to \infty$: basket-weighted CPI (headline).
        \item CV selects $\lambda$ on a log grid $[10^{-4}, 10^{3}]$.
      \end{itemize}
    \end{column}
  \end{columns}
\end{frame}

% ── Slide 10: Albacoreranks ────────────────────────────────────────────────
\begin{frame}{Model 2: Albacoreranks (Rank Space)}
  \begin{block}{Optimization Problem}
    \[
      \min_{\mathbf{w}} \;\; \frac{1}{T}\sum_{t=1}^{T}\bigl(y_t - \mathbf{o}_t'\mathbf{w}\bigr)^2
      \;+\; \lambda \sum_{r=1}^{K-1}(w_{r+1} - w_r)^2
      \qquad\text{[fused ridge]}
    \]
    \vspace{-6pt}
    \[
      \text{s.t.} \quad w_r \geq 0, \quad \bar{\mathbf{o}}'\mathbf{w} = \bar{y}
      \qquad\text{[mean constraint]}
    \]
  \end{block}
  \vspace{4pt}
  \begin{columns}[T]
    \begin{column}{0.48\textwidth}
      \textbf{Idea:}
      \begin{itemize}
        \item At each $t$, sort component growth rates low $\to$ high.
        \item $\mathbf{O}$: order-statistic matrix ($T \times K$).
        \item Learn which \emph{percentile} of the cross-sectional
              distribution predicts future inflation.
      \end{itemize}
    \end{column}
    \begin{column}{0.48\textwidth}
      \textbf{Differences from Albacorecomps:}
      \begin{itemize}
        \item \emph{Identity-free:} weights attach to ranks, not named components.
        \item Fused ridge: encourages smooth weights across adjacent ranks.
        \item Mean constraint replaces sum-to-one.
        \item $\lambda$ grid: $[10^{0}, 10^{5}]$ (penalizes non-smoothness).
      \end{itemize}
    \end{column}
  \end{columns}
\end{frame}

% ── Slide 11: Benchmarks ──────────────────────────────────────────────────
\begin{frame}{Benchmark Models}
  \begin{table}[h]
    \centering
    \small
    \begin{tabular}{lll}
      \toprule
      \textbf{Benchmark} & \textbf{Predictor at time $t$} & \textbf{Type} \\
      \midrule
      Random walk & Current headline 3m/3m & Na\"ive \\
      Core (ex fresh food) & Current core CPI 3m/3m & Na\"ive \\
      Core (ex food \& energy) & Current ``supercore'' 3m/3m & Na\"ive \\
      Unconditional mean & $\bar{y}_{1:t}$ (expanding mean) & Na\"ive \\
      OLS combination & $\hat{y}_t$ from OLS on 3 core measures & Parametric \\
      \bottomrule
    \end{tabular}
  \end{table}
  \vspace{6pt}
  \begin{itemize}
    \item \textbf{Unconditional mean} is the $R^2 = 0$ floor --- any model that
          cannot beat it is uninformative.
    \item \textbf{OLS benchmark} regresses $y$ on headline + core + supercore
          (expanding window, re-estimated at each step).
    \item All benchmarks use only information available at time $t$.
  \end{itemize}
\end{frame}

% ════════════════════════════════════════════════════════════════════════════
\section{Results}
% ════════════════════════════════════════════════════════════════════════════

% ── Slide 12: In-Sample Results ────────────────────────────────────────────
\begin{frame}{In-Sample Results}
  \begin{columns}[T]
    \begin{column}{0.48\textwidth}
      \textbf{Albacorecomps (Level 2, 47 components):}
      \begin{itemize}
        \item $\lambda$ selected via 10-fold expanding-window CV.
        \item Achieves tight in-sample fit with only a subset of non-zero weights.
        \item Weights deviate from prior mainly for
              \alert{energy}, \alert{food}, and \alert{services} components.
      \end{itemize}
      \vfill
      \textit{See: reg\_fig1\_weights.png, reg\_fig2\_lambda\_cv.png, reg\_fig3\_insample.png}
    \end{column}
    \begin{column}{0.48\textwidth}
      \textbf{Weight reallocation pattern:}
      \begin{itemize}
        \item Energy components (electricity, gas) receive
              \emph{increased} weight relative to prior --- they lead headline.
        \item Rent receives \emph{reduced} weight --- low signal, high basket share.
        \item Some small components (e.g.\ tobacco, bedding) shrink to zero.
      \end{itemize}
      \vspace{6pt}
      \textbf{Albacoreranks:}
      \begin{itemize}
        \item Smooth weight profile across ranks.
        \item Higher weight on \emph{central} and \emph{upper} ranks.
      \end{itemize}
    \end{column}
  \end{columns}
\end{frame}

% ── Slide 13: OOS Results ─────────────────────────────────────────────────
\begin{frame}{Out-of-Sample Performance}
  \begin{block}{Evaluation Metrics (OOS, 12-month forward headline)}
    \centering
    \small
    \begin{tabular}{lccc}
      \toprule
      \textbf{Model} & \textbf{RMSE} & \textbf{MAE} & \textbf{$R^2$} \\
      \midrule
      \textbf{Comps -- expanding} & \multicolumn{3}{c}{\textit{from scorecard}} \\
      \textbf{Comps -- rolling 20y} & \multicolumn{3}{c}{\textit{from scorecard}} \\
      \textbf{Ranks -- rolling 20y} & \multicolumn{3}{c}{\textit{from scorecard}} \\
      \midrule
      Random walk & \multicolumn{3}{c}{\textit{from scorecard}} \\
      Core (ex fresh food) & \multicolumn{3}{c}{\textit{from scorecard}} \\
      Core (ex food \& energy) & \multicolumn{3}{c}{\textit{from scorecard}} \\
      Unconditional mean & \multicolumn{3}{c}{\textit{from scorecard}} \\
      \bottomrule
    \end{tabular}
  \end{block}
  \vspace{4pt}
  \small
  \textbf{Instructions:} Replace \textit{from scorecard} with actual numbers from the
  \texttt{print\_scorecard} output after running the regressions.\\[4pt]
  \textit{See: reg\_fig4\_oos.png for time-series and rolling RMSE plots.}
\end{frame}

% ── Slide 14: OOS Discussion ──────────────────────────────────────────────
\begin{frame}{OOS Analysis: Time Variation and Rolling RMSE}
  \begin{columns}[T]
    \begin{column}{0.55\textwidth}
      \textbf{Key findings:}
      \begin{enumerate}
        \item Assemblage models \alert{outperform} na\"ive benchmarks in RMSE and $R^2$
              over the full OOS period.
        \item \textbf{Rolling 20-year window} adapts better than expanding window during
              regime changes (deflation $\to$ Abenomics $\to$ 2022 surge).
        \item \textbf{Rank-space model} provides robustness: it does not rely on
              component identity, reducing sensitivity to structural breaks.
        \item Na\"ive core measures lag during rapid inflation shifts.
      \end{enumerate}
    \end{column}
    \begin{column}{0.42\textwidth}
      \textbf{Rolling 36-month RMSE:}
      \begin{itemize}
        \item All models struggle during the 2022 shock.
        \item Assemblage recovers faster than random walk.
        \item Core (ex food \& energy) performs worst
              during energy-driven episodes.
      \end{itemize}
    \end{column}
  \end{columns}
\end{frame}

% ── Slide 15: Multi-Level Comparison ──────────────────────────────────────
\begin{frame}{Multi-Level Analysis: Level 1 vs.\ Level 2}
  \begin{columns}[T]
    \begin{column}{0.48\textwidth}
      \textbf{Level 1} (9 major sectors):
      \begin{itemize}
        \item Coarser grouping.
        \item Fewer parameters $\Rightarrow$ less overfitting risk.
        \item Weights are more interpretable (sector-level).
        \item Good baseline, but cannot capture within-sector heterogeneity.
      \end{itemize}
    \end{column}
    \begin{column}{0.48\textwidth}
      \textbf{Level 2} (47 mid-level components):
      \begin{itemize}
        \item Richer signal: can distinguish energy subtypes, food subtypes, etc.
        \item More parameters $\Rightarrow$ stronger regularization needed.
        \item Better OOS performance when $\lambda$ is well-tuned.
        \item Rank-space model benefits most from granularity.
      \end{itemize}
    \end{column}
  \end{columns}
  \vfill
  \begin{block}{Takeaway}
    The Level~2 component decomposition provides the best trade-off between
    signal richness and estimation stability for both Albacorecomps and Albacoreranks.
  \end{block}
\end{frame}

% ════════════════════════════════════════════════════════════════════════════
\section{Conclusion}
% ════════════════════════════════════════════════════════════════════════════

% ── Slide 16: Conclusion ──────────────────────────────────────────────────
\begin{frame}{Conclusion \& Future Work}
  \textbf{Summary:}
  \begin{enumerate}
    \item Assemblage regression learns \alert{adaptive core measures} that outperform
          fixed-weight published cores for forecasting 12-month forward headline
          CPI in Japan.
    \item Shrinkage toward basket-weight prior prevents overfitting while allowing
          data-driven reallocation.
    \item Rent isolation is critical --- its 18.3\% basket share with near-zero
          variance distorts unconstrained estimation.
    \item Rank-space approach (Albacoreranks) offers a complementary, identity-free
          perspective that is robust to compositional changes.
  \end{enumerate}

  \vspace{6pt}
  \textbf{Future directions:}
  \begin{itemize}
    \item Extend to Level~3 (744 items) with grouped-lasso penalties.
    \item Time-varying $\lambda$ via online CV.
    \item Cross-country comparison (US PCE, Euro HICP).
    \item Incorporate external regressors (exchange rate, oil price, output gap).
  \end{itemize}
\end{frame}

% ── Slide 17: Thank You ──────────────────────────────────────────────────
\begin{frame}[plain]
  \centering
  \vspace{2cm}
  {\Huge\bfseries\textcolor{darkblue}{Thank You}}\\[12pt]
  {\large Questions \& Discussion}\\[24pt]
  \textcolor{gray}{\small Code available in the project repository.}
\end{frame}

\end{document}
